\section{Introduction}\label{introduction}

One of the most puzzling questions in astrobiology is whether intelligent civilizations exist elsewhere in the galaxy, and if so, why we have not yet detected them. This apparent contradiction, often referred to as the Fermi Paradox or the Great Silence, raises the question of why, in a galaxy billions of years old, we observe no clear signs of advanced extraterrestrial life. Many solutions have been proposed to explain this Great Silence \citep{Webb2015-vi}. One class of explanations suggests that intelligent life is intrinsically rare, perhaps requiring an improbable convergence of planetary, chemical, or biological conditions \citep{Wesson1990-xi}. Another contends that while intelligence may emerge, technological civilizations are typically short-lived, succumbing to self-destruction through ecological collapse, conflict, or other systemic instabilities \citep{Vinn2024-pi}. A third class of explanations holds that civilizations may be neither rare nor short-lived, yet still remain undetected. Expansion and contact may be selective, directed preferentially toward planets that already display detectable technospheres \citep{Wandel2022-fp}, and even sustained interstellar settlement can naturally produce a sparsely or porously populated galaxy in which large regions, including our own, remain unvisited for long periods \citep{Carroll-Nellenback2019-aurora}. These classes are not mutually exclusive, and the mechanism we examine in this paper, intermittent technological activity, can operate in combination with any of them.

This paper focuses on the second class, which raises a question that is critical regardless of how common the emergence of intelligence may be: how long do technological civilizations last once they arise? Even if intelligent life is widespread, the number of civilizations detectable at any given time depends decisively on their longevity \citep{Balbi2021-bc}, and the observed silence itself has been used to place statistical bounds on that longevity \citep{Rahvar2026-lt}. This question is also central to projecting the long-term trajectory of our own civilization. Most treatments of longevity, however, assume that once a society becomes technologically active, it remains so continuously until it either collapses or deliberately goes silent. But this assumption is difficult to justify. Historical studies of Earth's civilizations reveal cycles of rise and fall with typical durations on the order of centuries \citep{Shermer2002-oi}, and even within a single civilization's lifespan, the temporal distribution of active periods may matter as much as their cumulative duration \citep{Balbi2018-fp}. Therefore, understanding civilizational longevity requires moving beyond static estimates and examining the temporal dynamics of technological activity itself. Throughout this paper we accordingly distinguish three related but distinct concepts: \emph{longevity}, the total lifespan of a civilization; \emph{technological activity}, the periods within that lifespan during which the civilization sustains a functioning technosphere; and \emph{detectability}, the production of remotely observable technosignatures, which requires technological activity but can also outlast it (Section~\ref{technosignatures-implications}).

There are good reasons to expect that civilizational collapse-recovery dynamics are complex. It has been argued that planetary civilizations face a fundamental bifurcation: either "asymptotic burnout" in which superlinear growth leads to singularity-driven collapse, or "homeostatic awakening", in which a civilization deliberately curbs expansion to achieve long-term equilibrium \citep{Wong2022-hv}. Sustainability constraints make exponential expansion highly unlikely, so advanced societies may stabilize or collapse before spreading across the Galaxy \citep{Haqq-Misra2009-jb,mullan2019population}. If collapse is common and recovery uncertain, then the technosphere (the totality of a civilization's technological infrastructure and output) may best be understood not as a permanent state but as an intermittent phenomenon, alternating between periods of activity and dormancy in a pattern shaped by resource availability, governance structure, and exposure to existential risk.

Technological intermittency can be quantified through a civilization's \emph{duty cycle}, defined as the fraction of its total lifespan during which it is technologically active. The duty cycle captures a reality that static longevity estimates miss---that a civilization persisting for a thousand years may spend substantial portions of that time in dormancy, recovery, or regression. In the context of the search for extraterrestrial intelligence, the relevance of intermittent activity is already recognized. Interstellar signals may be intermittent due to power constraints, planetary rotation, or targeted sequential transmissions \citep{Gray2020-cn}, and recent work applying Lindy's law to technosignature longevities suggests that technologically active periods are more likely to be short-lived than long-lived \citep{Balbi2024-bg}. But the implications of intermittency extend well beyond detectability. For our own civilization, understanding the dynamics of collapse and recovery, and how governance, resources, and resilience infrastructure shape the duty cycle, is directly relevant to long-term planning and existential risk mitigation. Both facets of the problem fall within the scope of astrobiology. Technosignature searches are now an established component of the field's observational agenda \citep{Haqq-Misra2025-rt}, and the detectability of technological life depends not only on where such life arises but on how continuously it remains active.

In this paper, we develop a framework for modeling these dynamics. Building on a recent scenario framework for Earth-originating civilization \citep{Haqq-Misra2025-rt}, we simulate collapse-recovery cycles across ten diverse sociotechnical trajectories spanning 1000 years and compute duty cycles for each. We then derive an effective detectability duration, $L_{\mathrm{eff}}$, and explore the implications for Earth's own civilizational resilience, for the Drake Equation's longevity term \emph{L}, and for the broader search for technosignatures. The ten scenarios are structured narrative projections of Earth-originating futures, not a statistical sample of civilizations, and the simulations built on them are exploratory. They are not empirical predictions about humanity's future or about extraterrestrial civilizations; their value lies in making the consequences of intermittency explicit and quantitative within a transparent and self-consistent set of assumptions. When we extend the results beyond Earth (Section~\ref{technosignatures-implications}), we do so conditionally. The specific duty-cycle values are Earth-derived illustrations, while the structural insight, that collapse-recovery dynamics generically reduce effective detectability, does not depend on the particular parameter values. We return to the limits of this extrapolation in Section~\ref{limitations}.
